\chapter{第四单元}

\Needspace{12\baselineskip}

\section{第四单元素养评价（A）}

\Needspace{8\baselineskip}

\pdfbookmark[2]{原书第25页}{orig-25}

\subsection*{原书第25页}

% p025_s01
\BookQuestion{一、1}
\(a\times(b+c)=\)（\hspace{1.8em}）\(\times\)（\hspace{1.8em}）\(+\)（\hspace{1.8em}）\(\times\)（\hspace{1.8em}）。\par
\begin{BookAnswer}
用乘法分配律：\[a\times(b+c)=a\times b+a\times c.\]四个空依次是\(a,b,a,c\)。括号外的\(a\)要分别与\(b,c\)相乘。
\end{BookAnswer}

% p025_s02
\BookQuestion{一、2}
观察右边的竖式，在计算过程中运用了（\hspace{1.8em}）律，用横式表示右边竖式的计算过程是（\hspace{1.8em}）。\par
\BookFigure{figures/p025_s02-01.tex}
\begin{BookAnswer}
竖式先算\(45\times2\)，再算\(45\times40\)，最后把两部分相加，运用了乘法分配律。\[\begin{aligned} &45\times42\\={}&45\times40+45\times2\\={}&1800+90\\={}&1890\end{aligned}\]
\end{BookAnswer}

% p025_s03
\BookQuestion{一、3}
小明在用计算器计算\(87\times99\)时，发现按键“9”出现故障，聪明的你想怎样用计算器算出结果？请写出你的算式：（\hspace{1.8em}）。\par
\begin{BookAnswer}
先把99转化成不用9键输入的等值计算。\par 输入法一：先乘100再减一份87\par \[\begin{aligned} &87\times(100-1)\\={}&87\times100-87\\={}&8700-87\\={}&8613\end{aligned}\]按键可依次为\(87,\times,100,-,87,=\)，全程不用9键。\par 输入法二：先算88加11，再乘87\par \[\begin{aligned} &(88+11)\times87\\={}&99\times87\\={}&8613\end{aligned}\]先按\(88,+,11,=\)，显示99；再按\(\times,87,=\)。99是计算结果，不是用9键输入，因此没有括号键的普通计算器也能这样算。\par 输入法三：把99分解成33乘3。由乘法结合律，可先算87乘33，再乘3。\[\begin{aligned} &87\times99\\={}&87\times(33\times3)\\={}&87\times33\times3\\={}&2871\times3\\={}&8613\end{aligned}\]按键可依次为\(87,\times,33,\times,3,=\)，所有输入数字中都没有9。
\end{BookAnswer}

% p025_s04
\BookQuestion{一、4（左上）}
在方框里填上合适的数或字母，在圆圈里填上“\(+\)、\(-\)、\(\times\)、\(\div\)”：\(98+74+26=\square+(74+\square)\)。\par
\begin{BookAnswer}
只改变结合顺序，没有交换加数：\(98+74+26=98+(74+26)\)。因此这里用加法结合律；两空依次填98、26。
\end{BookAnswer}

% p025_s05
\BookQuestion{一、4（左中）}
在方框里填上合适的数或字母，在圆圈里填运算符号：\(m\times21\times9=m\times(\square\bigcirc\square)\)。\par
\begin{BookAnswer}
乘法结合律：\(m\times21\times9=m\times(21\times9)\)。依次填21、乘号、9。
\end{BookAnswer}

% p025_s06
\BookQuestion{一、4（左下）}
在方框里填上合适的数或字母，在圆圈里填运算符号：\(380\div5\div2=380\bigcirc(\square\bigcirc\square)\)。\par
\begin{BookAnswer}
连续除以两个数，等于除以它们的积：\(380\div5\div2=380\div(5\times2)\)。依次填除号、5、乘号、2。
\end{BookAnswer}

% p025_s07
\BookQuestion{一、4（右上）}
在方框里填上合适的数或字母，在圆圈里填运算符号：\(612-124-76=612\bigcirc(\square\bigcirc\square)\)。\par
\begin{BookAnswer}
减法的性质：\(612-124-76=612-(124+76)\)。依次填减号、124、加号、76。
\end{BookAnswer}

% p025_s08
\BookQuestion{一、4（右中）}
在方框里填上合适的数或字母，在圆圈里填运算符号：\(389+102=389\bigcirc\square\bigcirc\square\)。\par
\begin{BookAnswer}
把102拆成100与2：\(389+102=389+100+2\)。依次填加号、100、加号、2。
\end{BookAnswer}

% p025_s09
\BookQuestion{一、4（右下）}
在方框里填上合适的数或字母：\((a+b)\times9=a\times\square+\square\times\square\)。\par
\begin{BookAnswer}
乘法分配律：\((a+b)\times9=a\times9+b\times9\)。三空依次填9、\(b\)、9。
\end{BookAnswer}

% p025_s10
\BookQuestion{一、5（左上）}
在圆圈里填上“\(>\)”“\(<\)”或“\(=\)”：\(50\times(12+33)\ \bigcirc\ 50\times12+33\)。\par
\begin{BookAnswer}
填\(>\)。左边为\(50\times45=2250\)，右边为\(600+33=633\)；左边的33也乘了50。
\end{BookAnswer}

% p025_s11
\BookQuestion{一、5（右上）}
在圆圈里填上“\(>\)”“\(<\)”或“\(=\)”：\(132-68-12\ \bigcirc\ 132-68+12\)。\par
\begin{BookAnswer}
填\(<\)。左边\(132-68-12=52\)，右边\(132-68+12=76\)。
\end{BookAnswer}

% p025_s12
\BookQuestion{一、5（左中）}
在圆圈里填上“\(>\)”“\(<\)”或“\(=\)”：\(720\div8\times6\ \bigcirc\ 720\div8\div6\)。\par
\begin{BookAnswer}
填\(>\)。左边先除8再乘6，得540；右边先除8再除6，得15。
\end{BookAnswer}

% p025_s13
\BookQuestion{一、5（右中）}
在圆圈里填上“\(>\)”“\(<\)”或“\(=\)”：\(32\times20\times8\ \bigcirc\ 32\times(20\times8)\)。\par
\begin{BookAnswer}
填\(=\)。\(32\times20\times8=32\times(20\times8)\)，依据乘法结合律。
\end{BookAnswer}

% p025_s14
\BookQuestion{一、5（左下）}
在圆圈里填上“\(>\)”“\(<\)”或“\(=\)”：\(12\times56\ \bigcirc\ 56\times12\)。\par
\begin{BookAnswer}
填\(=\)。\(12\times56=56\times12\)，依据乘法交换律。
\end{BookAnswer}

% p025_s15
\BookQuestion{一、5（右下）}
在圆圈里填上“\(>\)”“\(<\)”或“\(=\)”：\(85\times103\ \bigcirc\ 85\times100+3\)。\par
\begin{BookAnswer}
填\(>\)。左边\(85\times103=8500+255\)，右边为\(8500+3\)，而\(255>3\)。
\end{BookAnswer}

% p025_s16
\BookQuestion{一、6}
如果\(\bigcirc+\triangle=4\)，那么\(25\times\bigcirc+25\times\triangle=\)（\hspace{1.8em}）。\par
\begin{BookAnswer}
反用乘法分配律：\[\begin{aligned} &25\times\bigcirc+25\times\triangle\\={}&25\times(\bigcirc+\triangle)\\={}&25\times4\\={}&100\end{aligned}\]
\end{BookAnswer}

% p025_s17
\BookQuestion{一、7}
明明在计算时，把\(8\times(\square+6)\)错算成\(8\times\square+6\)，他得到的结果与正确结果相差（\hspace{1.8em}）。\par
\begin{BookAnswer}
正确式中的6应乘8，错式只有一份6；两结果相差7份6：\[\begin{aligned}&8\times6-6\\={}&6\times(8-1)\\={}&6\times7\\={}&42\end{aligned}\]
\end{BookAnswer}

\Needspace{8\baselineskip}

\pdfbookmark[2]{原书第26页}{orig-26}

\subsection*{原书第26页}

% p026_s01
\BookQuestion{二、1}
计算\(32\times125\)，下面（\hspace{1.8em}）中计算方法是错误的。\par
\noindent A.\(125\times8+125\times4\)\par
\noindent B.\(125\times8\times4\)\par
\noindent C.\(125\times30+125\times2\)\par
\noindent D.\(125\times40-125\times8\)\par
\begin{BookAnswer}
选\(A\)。A把32误拆成了\(8+4\)，实际只算了12个125。其余三式分别利用\(32=8\times4\)、\(32=30+2\)、\(32=40-8\)。
\end{BookAnswer}

% p026_s02
\BookQuestion{二、2}
下面的算式中，正确的是（\hspace{1.8em}）。\par
\noindent A.\(54\times13+7=54\times(13+7)\)\par
\noindent B.\(99\times16=100\times16-1\)\par
\noindent C.\(350\div2\div5=350\div(2\times5)\)\par
\noindent D.\(45\times36=63\times54\)\par
\begin{BookAnswer}
选\(C\)。连续除以2再除以5，等于除以\(2\times5\)：\[\begin{aligned} &350\div2\div5\\={}&350\div(2\times5)\\={}&350\div10\\={}&35\end{aligned}\]
\end{BookAnswer}

% p026_s03
\BookQuestion{二、3}
算式“\(783-256-\square\)”可以用减法的性质进行简算，则方框表示的数是（\hspace{1.8em}）。\par
\noindent A.117\par
\noindent B.256\par
\noindent C.283\par
\noindent D.244\par
\begin{BookAnswer}
\textbf{C、D都能简便计算，本题答案不唯一。}\par 选\(D\)时，先把两个减数相加凑整，运用减法的性质：\[\begin{aligned}&783-256-244\\&=783-(256+244)\\&=783-500\\&=283\end{aligned}\]\par 选\(C\)时，交换两个减数的位置，先减去\(283\)凑整，运用减法的性质：\[\begin{aligned}&783-256-283\\&=783-283-256\\&=500-256\\&=244\end{aligned}\]\par 两种方法都先凑出\(500\)，都符合小学简便计算要求。题干没有限定必须先把两个减数相加，因此不能只选\(D\)。
\end{BookAnswer}

% p026_s04
\BookQuestion{二、4}
以下选项（\hspace{1.8em}）的画图方式，可以说明\((2+3)\times m=2\times m+3\times m\)。\par
\noindent A.线段图\par
\noindent B.分割长方形面积图\par
\noindent C.文具总价图\par
\noindent D.两个长方形宽均m、高5的图\par
\BookFigure{figures/p026_s04-01.tex}
\begin{BookAnswer}
选\(B\)。B图长方形的长为\(2+3\)，宽为\(m\)。整体面积是\((2+3)\times m\)；分成两块，面积为\(2\times m+3\times m\)。
\end{BookAnswer}

% p026_s05
\BookQuestion{二、5}
在古埃及，人们采用“双倍法”计算乘法。如计算\(23\times14\)的过程：\(23\times2=46\)，\allowbreak{}\(23\times4=92\)，\allowbreak{}\(23\times8=184\)，\allowbreak{}\(23\times14=46+92+184=322\)。实际上“双倍法”计算使用的运算律是（\hspace{1.8em}）。\par
\noindent A.乘法结合律\par
\noindent B.乘法交换律\par
\noindent C.乘法分配律\par
\noindent D.乘法结合律和分配律\par
\begin{BookAnswer}
选\(C\)。因为\(14=2+4+8\)，所以\(23\times14=23\times2+23\times4+23\times8\)。把一个数与和相乘改为分别相乘后相加，依据乘法分配律。
\end{BookAnswer}

% p026_s06
\BookQuestion{二、6}
如右图，要计算这个图形的面积，下面算式中正确的有（\hspace{1.8em}）个。①\(15\times32+15\times18\)；②\(15\times32+18\)；③\(15\times(32+18)\)；④\((15+18)\times32-18\times(32-15)\)。\par
\noindent A.1\par
\noindent B.2\par
\noindent C.3\par
\noindent D.4\par
\BookFigure{figures/p026_s06-01.tex}
\begin{BookAnswer}
选\(C\)。正确的是①③④，共3个。①把图形分成两个长方形；③把①按分配律合并；④先补成长方形再减去缺角。②把面积和长度相加，不能表示该图形面积。三种正确算法都得750平方米。
\end{BookAnswer}

% p026_s07
\BookQuestion{三、1（1）}
直接写得数：\(45+55=\)。\par
\begin{BookAnswer}
两数相加凑成100。\[\begin{aligned} &45+55\\={}&100\end{aligned}\]
\end{BookAnswer}

% p026_s08
\BookQuestion{三、1（2）}
直接写得数：\(125\times8=\)。\par
\begin{BookAnswer}
125与8凑1000。\[\begin{aligned} &125\times8\\={}&1000\end{aligned}\]
\end{BookAnswer}

% p026_s09
\BookQuestion{三、1（3）}
直接写得数：\(5\times7\times4=\)。\par
\begin{BookAnswer}
先把5与4相乘得20。\[\begin{aligned} &5\times7\times4\\={}&140\end{aligned}\]
\end{BookAnswer}

% p026_s10
\BookQuestion{三、1（4）}
直接写得数：\(25\times6\times4=\)。\par
\begin{BookAnswer}
先把25与4相乘得100。\[\begin{aligned} &25\times6\times4\\={}&600\end{aligned}\]
\end{BookAnswer}

% p026_s11
\BookQuestion{三、1（5）}
直接写得数：\(4\times25=\)。\par
\begin{BookAnswer}
4与25凑100。\[\begin{aligned} &4\times25\\={}&100\end{aligned}\]
\end{BookAnswer}

% p026_s12
\BookQuestion{三、1（6）}
直接写得数：\(123-97=\)。\par
\begin{BookAnswer}
减97可先减100再加3。\[\begin{aligned} &123-97\\={}&26\end{aligned}\]
\end{BookAnswer}

% p026_s13
\BookQuestion{三、1（7）}
直接写得数：\(6\times4\div4\times6=\)。\par
\begin{BookAnswer}
同级运算从左往右，乘4再除4相互抵消。\[\begin{aligned} &6\times4\div4\times6\\={}&36\end{aligned}\]
\end{BookAnswer}

% p026_s14
\BookQuestion{三、1（8）}
直接写得数：\(130-30\div5=\)。\par
\begin{BookAnswer}
先除后减。\[\begin{aligned} &130-30\div5\\={}&124\end{aligned}\]
\end{BookAnswer}

\Needspace{8\baselineskip}

\pdfbookmark[2]{原书第27页}{orig-27}

\subsection*{原书第27页}

% p027_s01
\BookQuestion{三、2（1）}
用简便方法计算下面各题：\(435-142-58\)。\par
\begin{BookAnswer}
先把减数合并为整百数。\[\begin{aligned} &435-142-58\\={}&435-(142+58)\\={}&435-200\\={}&235\end{aligned}\]
\end{BookAnswer}

% p027_s02
\BookQuestion{三、2（2）}
用简便方法计算下面各题：\(68\times67+32\times67\)。\par
\begin{BookAnswer}
反用乘法分配律。\[\begin{aligned} &68\times67+32\times67\\={}&(68+32)\times67\\={}&100\times67\\={}&6700\end{aligned}\]
\end{BookAnswer}

% p027_s03
\BookQuestion{三、2（3）}
用简便方法计算下面各题：\(215-98\)。\par
\begin{BookAnswer}
把98看作100，减多的2要补回。\[\begin{aligned} &215-98\\={}&215-100+2\\={}&115+2\\={}&117\end{aligned}\]
\end{BookAnswer}

% p027_s04
\BookQuestion{三、2（4）}
用简便方法计算下面各题：\(99\times99+99\)。\par
\begin{BookAnswer}
把末尾99看作99乘1。\[\begin{aligned} &99\times99+99\\={}&99\times(99+1)\\={}&99\times100\\={}&9900\end{aligned}\]
\end{BookAnswer}

% p027_s05
\BookQuestion{三、2（5）}
用简便方法计算下面各题：\(3600\div25\div4\)。\par
\begin{BookAnswer}
连续除以两数，改成除以它们的积。\[\begin{aligned} &3600\div25\div4\\={}&3600\div(25\times4)\\={}&3600\div100\\={}&36\end{aligned}\]
\end{BookAnswer}

% p027_s06
\BookQuestion{三、2（6）}
用简便方法计算下面各题：\(25\times32\times125\)。\par
\begin{BookAnswer}
把\(32\)分成\(4\times8\)，用乘法结合律把\(25\)与\(4\)、\(8\)与\(125\)分别结合凑整。\par \[\begin{aligned} &25\times32\times125\\={}&(25\times4)\times(8\times125)\\={}&100\times1000\\={}&100000\end{aligned}\]
\end{BookAnswer}

% p027_s07
\BookQuestion{三、2（7）}
用简便方法计算下面各题：\((25+11)\times40\)。\par
\begin{BookAnswer}
用乘法分配律。\[\begin{aligned} &(25+11)\times40\\={}&25\times40+11\times40\\={}&1000+440\\={}&1440\end{aligned}\]
\end{BookAnswer}

% p027_s08
\BookQuestion{三、2（8）}
用简便方法计算下面各题：\(357-156-144+43\)。\par
\begin{BookAnswer}
先用减法的性质合并减数156和144，再用加法交换律、结合律，让357与43凑400；每个减数仍须被减去。\par 把加数凑400、减数凑300。\[\begin{aligned} &357-156-144+43\\={}&(357+43)-(156+144)\\={}&400-300\\={}&100\end{aligned}\]
\end{BookAnswer}

% p027_s09
\BookQuestion{三、2（9）}
用简便方法计算下面各题：\(280\div35\)。\par
\begin{BookAnswer}
35拆成7乘5。\[\begin{aligned} &280\div35\\={}&280\div7\div5\\={}&40\div5\\={}&8\end{aligned}\]另一算法：\[\begin{aligned} &280\div35\\={}&(280\times2)\div(35\times2)\\={}&560\div70\\={}&8\end{aligned}\]
\end{BookAnswer}

% p027_s10
\BookQuestion{四、1}
为了迎接学校的体质健康达标测试，小明从周一到周四坚持每天练习跳绳，每天的跳绳次数如下：156下、195下、144下、205下。这四天他一共跳了多少下？\par
\begin{BookAnswer}
四天总数就是四天各自次数的和。先用运算律使计算简便，再按原顺序复核；这是同一个数量关系的两种计算次序。\par 主解：用交换律和结合律先凑整\par \[\begin{aligned} &156+195+144+205\\={}&(156+144)+(195+205)\\={}&300+400\\={}&700\text{(下)}\end{aligned}\]答：这四天一共跳了700下。\par 验算：按日期依次求累计数\par \[\begin{aligned} &156+195+144+205\\={}&351+144+205\\={}&495+205\\={}&700\text{(下)}\end{aligned}\]答：这四天一共跳了700下。\par
\end{BookAnswer}

% p027_s11
\BookQuestion{四、2}
AI智能物流仓库采用机器人分拣快递，仓库一共有12个分拣区，每个分拣区域分成4个工作小组，每个工作小组有25台智能分拣机器人。这个仓库一共有多少台分拣机器人？\par
\begin{BookAnswer}
方法一：先求每区台数，再求全仓库台数。\[\begin{aligned}&25\times4\times12\\={}&100\times12\\={}&1200\text{(台)}\end{aligned}\]答：这个仓库一共有1200台分拣机器人。\par 方法二：按全仓库小组数列式，再用乘法结合律先凑100。\[\begin{aligned}&25\times(4\times12)\\={}&(25\times4)\times12\\={}&100\times12\\={}&1200\text{(台)}\end{aligned}\]答：这个仓库一共有1200台分拣机器人。两者是数量分析路径不同，变形后的计算方法相同。
\end{BookAnswer}

\Needspace{8\baselineskip}

\pdfbookmark[2]{原书第28页}{orig-28}

\subsection*{原书第28页}

% p028_s01
\BookQuestion{四、3}
学校组织爱心义卖活动，四年级同学一共筹得1500元善款。他们先为山区小学购买学习用品花了665元，又为留守儿童购买书籍花了135元。还剩下多少善款？（5分）\par
\begin{BookAnswer}
先把两笔支出合计，凑成800元。\[\begin{aligned}&1500-(665+135)\\={}&1500-800\\={}&700\text{(元)}\end{aligned}\]答：还剩700元善款。\par 按购买顺序也可列成1500减665再减135；计算时同样合并减数：\[\begin{aligned}&1500-665-135\\={}&1500-(665+135)\\={}&1500-800\\={}&700\text{(元)}\end{aligned}\]这只是同一数量关系的两种列式，不另算一种独立解法。
\end{BookAnswer}

% p028_s02
\BookQuestion{四、4}
为了在实践活动中培养“有理想、有道德、有文化、有纪律”的四有新人，光明小学组织全校少先队员开展“红领巾志愿服务”活动。学校共25个中队，每个中队分成4个小队，全校1200名少先队员全部参与活动。平均每个小队有多少名少先队员？（5分）\par
\begin{BookAnswer}
方法一：先求全校小队数，再平均分。\[\begin{aligned}&1200\div(4\times25)\\={}&1200\div100\\={}&12\text{(名)}\end{aligned}\]答：平均每个小队有12名少先队员。\par 方法二：先分中队，再分小队；计算时用连除性质合并除数凑100。\[\begin{aligned}&1200\div25\div4\\={}&1200\div(25\times4)\\={}&1200\div100\\={}&12\text{(名)}\end{aligned}\]答：平均每个小队有12名少先队员。
\end{BookAnswer}

% p028_s03
\BookQuestion{四、5（1）}
为了提升乡村物流效率，某乡镇要新修一条快递专用配送通道，由甲、乙两个施工队从两端同时施工。甲队每天修53米，乙队每天修47米，两队合作16天顺利完工。（1）这条配送通道全长多少米？（5分）\par
\begin{BookAnswer}
方法一：先求两队每天共修多少米。\[\begin{aligned}&(53+47)\times16\\={}&100\times16\\={}&1600\text{(米)}\end{aligned}\]答：配送通道全长1600米。\par 方法二：从两队路程相加列式，反用乘法分配律凑100。\[\begin{aligned}&53\times16+47\times16\\={}&(53+47)\times16\\={}&100\times16\\={}&1600\text{(米)}\end{aligned}\]答：配送通道全长1600米。两种数量分析路径在简算时归为同一计算结构。
\end{BookAnswer}

% p028_s04
\BookQuestion{四、5（2）}
为了提升乡村物流效率，某乡镇要新修一条快递专用配送通道，由甲、乙两个施工队从两端同时施工。甲队每天修53米，乙队每天修47米，两队合作16天顺利完工。（2）施工结束时，甲队比乙队多修了多少米？（5分）\par
\begin{BookAnswer}
方法一：每天多修的长度乘天数。\[\begin{aligned}&(53-47)\times16\\={}&6\times16\\={}&96\text{(米)}\end{aligned}\]答：甲队比乙队多修96米。\par 方法二：从两队总路程相减列式，反用乘法分配律简算。\[\begin{aligned}&53\times16-47\times16\\={}&(53-47)\times16\\={}&6\times16\\={}&96\text{(米)}\end{aligned}\]答：甲队比乙队多修96米。两种数量分析路径在简算时归为同一计算结构。
\end{BookAnswer}

\Needspace{12\baselineskip}

\section{第四单元素养评价（B）}

\Needspace{8\baselineskip}

\pdfbookmark[2]{原书第29页}{orig-29}

\subsection*{原书第29页}

% p029_s01
\BookQuestion{一、1}
超市运来3箱文具，第一箱有277件，第二箱有164件，第三箱有536件。计算一共有多少件文具时，先把第（\hspace{1.8em}）箱和第（\hspace{1.8em}）箱的数量合起来，再加上第（\hspace{1.8em}）箱的数量，这样计算比较简便，列式为（\hspace{1.8em}）。\par
\begin{BookAnswer}
第二、三箱可以先凑整：\(164+536=700\)。再加第一箱277件。综合算式为：\[\begin{aligned} &277+(164+536)\\={}&277+700\\={}&977\text{(件)}\end{aligned}\]
\end{BookAnswer}

% p029_s02
\BookQuestion{一、2}
计算\(125\times5\times8\times20\)时，可以先算（\hspace{1.8em}）和（\hspace{1.8em}）（填算式），再算（\hspace{1.8em}）（填算式），这是运用了（\hspace{1.8em}）律和（\hspace{1.8em}）律使计算简便。\par
\begin{BookAnswer}
先算\(125\times8=1000\)和\(5\times20=100\)，再算\(1000\times100\)。交换因数位置，用乘法交换律；改变结合顺序，用乘法结合律。
\end{BookAnswer}

% p029_s03
\BookQuestion{一、3（左上）}
在圆圈里填上“＞”“＜”或“＝”：\(260\times48\ \bigcirc\ 26\times480\)。\par
\begin{BookAnswer}
填\(=\)。把260变为26、48变为480，一个因数缩小10倍，另一个扩大10倍，积不变，均为12480。
\end{BookAnswer}

% p029_s04
\BookQuestion{一、3（右上）}
在圆圈里填上“＞”“＜”或“＝”：\(301\times24\ \bigcirc\ 300\times24+1\)。\par
\begin{BookAnswer}
填\(>\)。\(301\times24=300\times24+24\)，比\(300\times24+1\)多23。
\end{BookAnswer}

% p029_s05
\BookQuestion{一、3（左中）}
在圆圈里填上“＞”“＜”或“＝”：\(145-18-82\ \bigcirc\ 145-(18+82)\)。\par
\begin{BookAnswer}
填\(=\)。\(145-18-82=145-(18+82)=45\)，依据减法的性质。
\end{BookAnswer}

% p029_s06
\BookQuestion{一、3（右中）}
在圆圈里填上“＞”“＜”或“＝”：\(50\times(33-13)\ \bigcirc\ 50\times33-13\)。\par
\begin{BookAnswer}
填\(<\)。左边\(50\times(33-13)=1000\)，右边\(50\times33-13=1637\)。左边减的是13个50，右边只减13。
\end{BookAnswer}

% p029_s07
\BookQuestion{一、3（左下）}
在圆圈里填上“＞”“＜”或“＝”：\((13+4)\times6\ \bigcirc\ 13\times6+4\times6\)。\par
\begin{BookAnswer}
填\(=\)。\((13+4)\times6=13\times6+4\times6\)，用乘法分配律，两边均为102。
\end{BookAnswer}

% p029_s08
\BookQuestion{一、3（右下）}
在圆圈里填上“＞”“＜”或“＝”：\(50\div(2\times5)\ \bigcirc\ 50\div2\times5\)。\par
\begin{BookAnswer}
填\(<\)。左边\(50\div(2\times5)=5\)，右边\(50\div2\times5=125\)。除以积应写成连续除，不能改成乘。
\end{BookAnswer}

% p029_s09
\BookQuestion{一、4}
实验小学运动会开幕式有8个方队，每个方队25人，每人拿2个气球，一共需要（\hspace{1.8em}）个气球。\par
\begin{BookAnswer}
每人2个气球，每方队25人，共8个方队：\[\begin{aligned} &2\times25\times8\\={}&50\times8\\={}&400\text{(个)}\end{aligned}\]还可先求总人数\(25\times8=200\text{(人)}\)，再求\(2\times200=400\text{(个)}\)。两条路径分别先求“每方队所需气球数”和“全体人数”，都遵循每份数×份数。
\end{BookAnswer}

% p029_s10
\BookQuestion{一、5}
一块菜地每行能种18棵白菜，一共有98行，这块菜地共能种多少棵白菜？小明列式\(18\times98\)，结合右图，还可列综合算式（\hspace{1.8em}），这是运用（\hspace{1.8em}）律。\par
\BookFigure{figures/p029_s10-01.tex}
\begin{BookAnswer}
把98行看作100行再去掉2行，运用乘法分配律：\[\begin{aligned} &18\times(100-2)\\={}&18\times100-18\times2\\={}&1800-36\\={}&1764\text{(棵)}\end{aligned}\]
\end{BookAnswer}

% p029_s11
\BookQuestion{一、6（1）上半}
明明用乘法分配律比较\(181\times67\)和\(182\times66\)的大小。（1）请帮他把思考和比较过程填完整。\(181\times67=181\times(\square+1)=181\times66+\square\times1\)。\par
\begin{BookAnswer}
把67拆成66加1：\(181\times67=181\times(66+1)=181\times66+181\times1\)。两空填66、181。
\end{BookAnswer}

\Needspace{8\baselineskip}

\pdfbookmark[2]{原书第30页}{orig-30}

\subsection*{原书第30页}

% p030_s01
\BookQuestion{一、6（1）下半}
（续）\(182\times66=(\square+1)\times66=181\times66+1\times\square\)。\par
\begin{BookAnswer}
把182拆成181加1：\(182\times66=(181+1)\times66=181\times66+1\times66\)。两空填181、66。
\end{BookAnswer}

% p030_s02
\BookQuestion{一、6（1）比较及规律}
所以\(181\times67\ \bigcirc\ 182\times66\)。发现：当两个数之和相同时，如果两个数之间的差越小，那么这两个数的积就越（\hspace{1.8em}）。\par
\begin{BookAnswer}
两式都有\(181\times66\)这一部分；左式另加181，右式另加66，所以填\(>\)，左边多\(181-66=115\)。对于和相同的两个正数，两个数之间的差越小，积越大。
\end{BookAnswer}

% p030_s03
\BookQuestion{一、6（2）}
（2）根据发现来比较：\(9999\times2026\ \bigcirc\ 9998\times2027\)。\par
\begin{BookAnswer}
填\(<\)。两组因数的和都为12025，右式两因数的差更小。也可比较共同部分：左边等于\(9998\times2026+2026\)，右边等于\(9998\times2026+9998\)；右边比左边多7972。
\end{BookAnswer}

% p030_s04
\BookQuestion{一、7（例式）}
\((81+36)\div9=81\div9+36\div9\)，\allowbreak{}\((121+22)\div11=121\div11+22\div11\)，请你照样子写一组这样的算式（\hspace{1.8em}）。\par
\begin{BookAnswer}
例如\((6+3)\div3=6\div3+3\div3\)。左边先求和再除3；右边分别除以3再相加，结果都是3。
\end{BookAnswer}

% p030_s05
\BookQuestion{一、7（字母规律）}
用字母表示上面的规律（\hspace{1.8em}）。\par
\begin{BookAnswer}
\[(a+b)\div c=a\div c+b\div c\quad(c\ne0).\]共同的除数必须不为0。这里是和作被除数，不能把除数的和随意拆开。
\end{BookAnswer}

% p030_s06
\BookQuestion{二、1}
\(53\times99=53\times100-53\times1\)应用了（\hspace{1.8em}）。\par
\noindent A. 加法交换律\par
\noindent B. 乘法交换律\par
\noindent C. 乘法结合律\par
\noindent D. 乘法分配律\par
\begin{BookAnswer}
选\(D\)。\(99=100-1\)，所以\(53\times99=53\times(100-1)=53\times100-53\times1\)，用了乘法分配律。
\end{BookAnswer}

% p030_s07
\BookQuestion{二、2}
下面（\hspace{1.8em}）的算式和\(27\times99\)的结果相等。\par
\noindent A. \(27\times90\times9\)\par
\noindent B. \(27\times100-1\)\par
\noindent C. \(27\times100-27\)\par
\noindent D. \(27\times100-100\)\par
\begin{BookAnswer}
选\(C\)。\(27\times99=27\times(100-1)=27\times100-27\)。少乘1个27，应减27。
\end{BookAnswer}

% p030_s08
\BookQuestion{二、3}
如果\(a\times(b\times18)=216\)，那么\(a\times b=\)（\hspace{1.8em}）。\par
\noindent A. 17\par
\noindent B. 198\par
\noindent C. 3888\par
\noindent D. 12\par
\begin{BookAnswer}
选\(D\)。由乘法结合律，\allowbreak{}\((a\times b)\times18=216\)，故\(a\times b=216\div18=12\)。
\end{BookAnswer}

% p030_s09
\BookQuestion{二、4}
小虎在一次计算中把\(5\times(a+6)\)错写成了\(5\times(a+9)\)，计算结果比原来（\hspace{1.8em}）。\par
\noindent A. 增加15\par
\noindent B. 减少15\par
\noindent C. 增加3\par
\noindent D. 减少3\par
\begin{BookAnswer}
选\(A\)。第二个加数由6变为9，多了3；先与5相乘，结果增加\(5\times(9-6)=15\)。
\end{BookAnswer}

% p030_s10
\BookQuestion{二、5}
如图，能用“\((15+24)\times32\)”解决的问题是（\hspace{1.8em}）。\par
\noindent A. 种西红柿、白菜和豆角的面积一共多少平方米\par
\noindent B. 种白菜和豆角的面积一共多少平方米\par
\noindent C. 种西红柿和豆角的面积一共多少平方米\par
\noindent D. 种西红柿和白菜的面积一共多少平方米\par
\BookFigure{figures/p030_s10-01.tex}
\begin{BookAnswer}
选D。15与24是两块区域的宽，32是共同的长，求两块合并面积。豆角区域宽20不在式中。继续计算时可用39=40−1补偿：\[\begin{aligned}&(15+24)\times32\\={}&39\times32\\={}&(40-1)\times32\\={}&40\times32-1\times32\\={}&1280-32\\={}&1248\text{(平方米)}\end{aligned}\]
\end{BookAnswer}

% p030_s11
\BookQuestion{二、6}
如图，能说明“\(6\times3+4\times3\)”与“\((6+4)\times3\)”相等的是（\hspace{1.8em}）。\par
\noindent A. ①②\par
\noindent B. ②③\par
\noindent C. ①②③\par
\noindent D. ①②③④\par
\BookFigure{figures/p030_s11-01.tex}
\BookFigure{figures/p030_s11-02.tex}
\begin{BookAnswer}
选\(B\)。②把长为\(6+4\)、宽为3的长方形分成两块；③每排有6个和4个圆，共3排。两图都能表示\(6\times3+4\times3=(6+4)\times3\)，均为30。①只是三段长度相加；④书和笔的数量并不都为3。
\end{BookAnswer}

\Needspace{8\baselineskip}

\pdfbookmark[2]{原书第31页}{orig-31}

\subsection*{原书第31页}

% p031_s01
\BookQuestion{三、1（1）}
直接写得数：\(28+172=\)。\par
\begin{BookAnswer}
两数凑整百。\[\begin{aligned} &28+172\\={}&200\end{aligned}\]
\end{BookAnswer}

% p031_s02
\BookQuestion{三、1（2）}
直接写得数：\(136-45-55=\)。\par
\begin{BookAnswer}
两个减数先凑100；原题得36。\[\begin{aligned} &136-45-55\\={}&136-(45+55)\\={}&136-100\\={}&36\end{aligned}\]
\end{BookAnswer}

% p031_s03
\BookQuestion{三、1（3）}
直接写得数：\(200\div25\div2=\)。\par
\begin{BookAnswer}
连续除可合并除数。\[\begin{aligned} &200\div25\div2\\={}&200\div(25\times2)\\={}&200\div50\\={}&4\end{aligned}\]
\end{BookAnswer}

% p031_s04
\BookQuestion{三、1（4）}
直接写得数：\(160\div2\times5=\)。\par
\begin{BookAnswer}
乘除同级，从左往右。\[\begin{aligned} &160\div2\times5\\={}&80\times5\\={}&400\end{aligned}\]
\end{BookAnswer}

% p031_s05
\BookQuestion{三、1（5）}
直接写得数：\(101\times12=\)。\par
\begin{BookAnswer}
把101拆成100和1。\[\begin{aligned} &101\times12\\={}&(100+1)\times12\\={}&1200+12\\={}&1212\end{aligned}\]
\end{BookAnswer}

% p031_s06
\BookQuestion{三、1（6）}
直接写得数：\(139+65-39=\)。\par
\begin{BookAnswer}
先让139减去39。\[\begin{aligned} &139+65-39\\={}&(139-39)+65\\={}&100+65\\={}&165\end{aligned}\]
\end{BookAnswer}

% p031_s07
\BookQuestion{三、1（7）}
直接写得数：\(37\times99+37=\)。\par
\begin{BookAnswer}
把后面的37看成37乘1。\[\begin{aligned} &37\times99+37\\={}&37\times(99+1)\\={}&37\times100\\={}&3700\end{aligned}\]
\end{BookAnswer}

% p031_s08
\BookQuestion{三、1（8）}
直接写得数：\(8\times5\div8\times5=\)。\par
\begin{BookAnswer}
先乘8又除8还原，两个5仍相乘：\[\begin{aligned}&8\times5\div8\times5\\={}&(8\div8)\times(5\times5)\\={}&1\times25\\={}&25\end{aligned}\]
\end{BookAnswer}

% p031_s09
\BookQuestion{三、2（1）}
用简便方法计算下面各题：（1）\(138+293+62+107\)。\par
\begin{BookAnswer}
交换、结合，分别凑整百。\[\begin{aligned} &138+293+62+107\\={}&(138+62)+(293+107)\\={}&200+400\\={}&600\end{aligned}\]
\end{BookAnswer}

% p031_s10
\BookQuestion{三、2（2）}
用简便方法计算下面各题：（2）\(574-189-211\)。\par
\begin{BookAnswer}
先合并两个减数。\[\begin{aligned} &574-189-211\\={}&574-(189+211)\\={}&574-400\\={}&174\end{aligned}\]
\end{BookAnswer}

% p031_s11
\BookQuestion{三、2（3）}
用简便方法计算下面各题：（3）\((25+7)\times8\)。\par
\begin{BookAnswer}
分配律使25乘8先凑200。\[\begin{aligned} &(25+7)\times8\\={}&25\times8+7\times8\\={}&200+56\\={}&256\end{aligned}\]
\end{BookAnswer}

% p031_s12
\BookQuestion{三、2（4）}
用简便方法计算下面各题：（4）\(250\times57\times40\)。\par
\begin{BookAnswer}
交换57与40的位置，用乘法交换律；先把250与40结合，用乘法结合律。\par 先把250与40凑成10000。\[\begin{aligned} &250\times57\times40\\={}&(250\times40)\times57\\={}&10000\times57\\={}&570000\end{aligned}\]
\end{BookAnswer}

% p031_s13
\BookQuestion{三、2（5）}
用简便方法计算下面各题：（5）\(720\div16\div5\)。\par
\begin{BookAnswer}
连续除以两数，等于除以它们的积。\[\begin{aligned} &720\div16\div5\\={}&720\div(16\times5)\\={}&720\div80\\={}&9\end{aligned}\]
\end{BookAnswer}

% p031_s14
\BookQuestion{三、2（6）}
用简便方法计算下面各题：（6）\(45\times102\)。\par
\begin{BookAnswer}
把102拆成100和2，用分配律。\[\begin{aligned} &45\times102\\={}&45\times(100+2)\\={}&45\times100+45\times2\\={}&4500+90\\={}&4590\end{aligned}\]
\end{BookAnswer}

% p031_s15
\BookQuestion{四、操作}
请你回忆学习“乘法分配律”时的探究过程，举例说明“\(12\times9+8\times9=(12+8)\times9\)”，可以用画图、列表、文字叙述等表达你的思考过程。（6分）\par
\begin{BookAnswer}
用植树情境解释原图：芍药每行12棵、种9行，牡丹每行8棵、也种9行。\par 分别数：先求两种花各有多少棵，再相加，列\(12\times9+8\times9\)。合并数：两种花每行合起来\(12+8\)棵，共9行，列\((12+8)\times9\)。\par 示意图的两块长方形分别对应\(12\times9\)和\(8\times9\)，高相同，合并后对应\((12+8)\times9\)；图中未逐个画出小格。分别求两块面积再相加，或先合并宽再求整体面积，结果相同：\[12\times9+8\times9=(12+8)\times9=180.\]这说明乘法分配律：和里的每个加数都要与括号外的9相乘。
\BookFigure{figures/p031_s15-01.tex}
\end{BookAnswer}

% p031_s16
\BookQuestion{五、1}
小军看一本故事书，第一天看了178页，第二天看了63页，还剩下122页没看，这本故事书一共有多少页？（5分）\par
\begin{BookAnswer}
全书页数＝第一天已看＋第二天已看＋还未看。先用运算律凑整，再按阅读进度复核。\par 主解：把178和122凑成300\par \[\begin{aligned} &178+63+122\\={}&(178+122)+63\\={}&300+63\\={}&363\text{(页)}\end{aligned}\]答：这本故事书一共有363页。\par 验算：先合计已看页数，再加未看页数\par \[\begin{aligned} &178+63+122\\={}&241+122\\={}&363\text{(页)}\end{aligned}\]答：这本故事书一共有363页。\par 回代检验：\(363-178-63=122\text{(页)}\)，与剩余页数相符。
\end{BookAnswer}

\Needspace{8\baselineskip}

\pdfbookmark[2]{原书第32页}{orig-32}

\subsection*{原书第32页}

% p032_s01
\BookQuestion{五、2}
2026年山东省足球超级联赛吉祥物“子竞”融合蹴鞠与现代足球元素。某赛事特许商店一天卖出“子竞”摆件和挂件各120个，每个摆件58元，每个挂件42元。这些摆件和挂件一共卖了多少元？（6分）\par
\begin{BookAnswer}
方法一：每种取1个配成一组，先求每组价格。\[\begin{aligned}&(58+42)\times120\\={}&100\times120\\={}&12000\text{(元)}\end{aligned}\]答：这些摆件和挂件一共卖了12000元。\par 方法二：分别计价列式，反用乘法分配律凑100。\[\begin{aligned}&58\times120+42\times120\\={}&(58+42)\times120\\={}&100\times120\\={}&12000\text{(元)}\end{aligned}\]答：这些摆件和挂件一共卖了12000元。
\end{BookAnswer}

% p032_s02
\BookQuestion{五、3（1）}
学校为课后延时服务的阅读社团采购图书，一共采购了12箱。每箱有25套书，每套有4本。（1）阅读社团一共采购了多少本图书？（5分）\par
\begin{BookAnswer}
方法一：先算每箱多少本\par \[\begin{aligned} &4\times25\times12\\={}&100\times12\\={}&1200\text{(本)}\end{aligned}\]答：一共采购了1200本图书。\par 方法二：先算共有多少套\par \[\begin{aligned}&4\times(25\times12)\\={}&(4\times25)\times12\\={}&100\times12\\={}&1200\text{(本)}\end{aligned}\]答：一共采购了1200本图书。\par
\end{BookAnswer}

% p032_s03
\BookQuestion{五、3（2）}
学校为课后延时服务的阅读社团采购图书，一共采购了12箱。每箱有25套书，每套有4本。（2）这些图书要分给阅读社团15个班级，每个班级有8个读书小组，平均每个读书小组分到多少本图书？（5分）\par
\begin{BookAnswer}
方法一：先求总小组数\par \[\begin{aligned} &1200\div(8\times15)\\={}&1200\div120\\={}&10\text{(本)}\end{aligned}\]答：平均每个读书小组分到10本图书。\par 方法二：先平均分到班，再分到组\par \[\begin{aligned} &1200\div15\div8\\={}&80\div8\\={}&10\text{(本)}\end{aligned}\]答：平均每个读书小组分到10本图书。\par
\end{BookAnswer}

% p032_s04
\BookQuestion{五、4}
学校“悦心农场”的一块劳动试验田如图所示，同学们计划在这里种植烟台特色的苹果树苗。这块劳动试验田的面积是多少平方米？（6分）\par
\BookFigure{figures/p032_s04-01.tex}
\begin{BookAnswer}
图中上方左块宽18米、高19米；下方横条长41米、高18米。\par 方法一：横向分成两个长方形，反用乘法分配律把41与19凑成60。\par \[\begin{aligned}&41\times18+18\times19\\={}&(41+19)\times18\\={}&60\times18\\={}&1080\text{(平方米)}\end{aligned}\]答：劳动试验田的面积是1080平方米。\par 方法二：纵向分割，两块的另一条边都为18米；先同时算19加18与41减18，再合并共同因数。\par \[\begin{aligned}&18\times(19+18)+(41-18)\times18\\={}&18\times37+23\times18\\={}&(37+23)\times18\\={}&60\times18\\={}&1080\text{(平方米)}\end{aligned}\]答：劳动试验田的面积是1080平方米。\par 方法三：补成大长方形后减缺角，展开后消去相同的41乘19，再凑60。\par \[\begin{aligned}&41\times(19+18)-(41-18)\times19\\={}&41\times19+41\times18-(41\times19-18\times19)\\={}&41\times18+18\times19\\={}&(41+19)\times18\\={}&60\times18\\={}&1080\text{(平方米)}\end{aligned}\]答：劳动试验田的面积是1080平方米。\par 这三种是不同的几何分割或补形路径；运算律变形不另算新的几何方法。
\end{BookAnswer}
